% Solutions to exercises for VI/04.
\begin{enumerate}[label=\textbf{E\arabic*.},leftmargin=*,itemsep=1em]
\item \textbf{\cref{ex:vi04-target-test}.}  For \(x\in X\), \(\varphi(x)\in Y=V(I)\) iff
\(g(f_1(x),\dots,f_m(x))=0\) for all \(g\in I\).  Requiring this for every \(x\) says
exactly that each substituted expression is zero in \(k[X]\).

\item \textbf{\cref{ex:vi04-pullback}.}  Composition gives
\((f+g)\circ\varphi=f\circ\varphi+g\circ\varphi\) and
\((fg)\circ\varphi=(f\circ\varphi)(g\circ\varphi)\); constants remain constant.

\item \textbf{\cref{ex:vi04-coordinates}.}  At \(x\),
\(\varphi^*(\bar y_i)(x)=\bar y_i(\varphi(x))=f_i(x)\).  The coordinate classes generate
\(k[Y]\), so their images determine the homomorphism.

\item \textbf{\cref{ex:vi04-contravariance}.}  For \(h\in k[Z]\),
\((\psi\circ\varphi)^*(h)=h\circ\psi\circ\varphi=\varphi^*(\psi^*h)\).

\item \textbf{\cref{ex:vi04-reconstruct}.}  Put \(f_i=\Phi(\bar y_i)\) and define
\(x\mapsto(f_1(x),\dots,f_m(x))\).  If \(g\in I(Y)\), then
\(g(\bar y)=0\) in \(k[Y]\), hence \(g(f)=\Phi(0)=0\); so the map lands in \(Y\).
Its pullback equals \(\Phi\) on generators and therefore everywhere.

\item \textbf{\cref{ex:vi04-continuity}.}  For \(W=V_Y(S)\),
\(x\in\varphi^{-1}(W)\) iff \(s(\varphi(x))=0\) for all \(s\in S\), iff
\((\varphi^*s)(x)=0\) for all \(s\).  Hence the inverse image is closed.

\item \textbf{\cref{ex:vi04-principal-open}.}  One has
\(x\in\varphi^{-1}(D_Y(f))\) iff \(f(\varphi(x))\neq0\), equivalently
\((\varphi^*f)(x)\neq0\).

\item \textbf{\cref{ex:vi04-parabola}.}  Let \(\alpha(t)=(t,t^2)\) and
\(\beta(x,y)=x\).  On \(y=x^2\), both composites are identities.  The pullbacks are
\(\alpha^*(\bar x)=t\), \(\alpha^*(\bar y)=t^2\), and \(\beta^*(t)=\bar x\).

\item \textbf{\cref{ex:vi04-axis}.}  The pullback sends \(x\mapsto t\), \(y\mapsto0\).
It is surjective and has kernel \((y)\), so \(k[t]\cong k[x,y]/(y)\).

\item \textbf{\cref{ex:vi04-affine-line-automorphism}.}  The inverse is
\(u\mapsto a^{-1}(u-b)\).  The inverse ring map sends \(t\mapsto a^{-1}(u-b)\).

\item \textbf{\cref{ex:vi04-squaring}.}  The pullback is \(u\mapsto t^2\) with image
\(k[t^2]\).  A nonzero polynomial in \(u\) stays nonzero after \(u=t^2\), while
\(t\notin k[t^2]\).

\item \textbf{\cref{ex:vi04-maps-to-a1}.}  A morphism to \(\mathbb A^1\) is a single
polynomial function \(f\in k[X]\), and every such function defines one.

\item \textbf{\cref{ex:vi04-iso-ring}.}  The inverse homomorphism
\((\varphi^*)^{-1}:k[X]\to k[Y]\) corresponds to \(\psi:Y\to X\).  Contravariance makes
both composites have identity pullback, hence they are identity morphisms.

\item \textbf{\cref{ex:vi04-closed-quotient}.}  Since \(Y\subseteq X\),
\(I(X)\subseteq I(Y)\), and quotienting gives
\(k[X]\twoheadrightarrow k[Y]\) with kernel \(I(Y)/I(X)=I_X(Y)\).

\item \textbf{\cref{ex:vi04-product-pair}.}  If \(f\) and \(g\) are morphisms, their
coordinate tuples concatenate to a morphism \((f,g)\).  Conversely, composing a map to
\(X\times Y\) with the projections recovers the pair.  These constructions are inverse.

\item \textbf{\cref{ex:vi04-constant-map}.}  For constant \(c_y\),
\(c_y^*(h)=h(y)\cdot1\).  Conversely, if every target coordinate pulls back to a scalar,
then every coordinate of \(\varphi\) is constant, so \(\varphi\) is constant.

\item \textbf{\cref{ex:vi04-graph}.}  In ambient coordinates \((x,y)\), the graph is
cut out inside \(X\times\mathbb A^m\) by \(y_i-f_i(x)=0\).  Projection to \(X\) is inverse
to the graph map, so it is a closed embedding.

\item \textbf{\cref{ex:vi04-kernel-image}.}  A function \(h\in k[Y]\) lies in the kernel
iff \(h(\varphi(x))=0\) for every \(x\).  Thus the kernel is exactly
\(I_Y(\varphi(X))\), and VI/02 gives the closure as its zero set.
\end{enumerate}
